\textbf{Phần 1. Trắc nghiệm lựa chọn}
\setcounterpageref{ex}{0}

\begin{ex} %[2X13N66]
Trong một hộp chỉ có các viên bi màu đỏ. Nhắm mắt lấy ra 1 viên bi, sự kiện "Lấy được viên bi màu đỏ" là
\choice
{Có thể}
{\True Chắc chắn}
{Không thể}
{Chắc chắn không xảy ra}
\loigiai{
Trong hộp chỉ có bi đỏ nên chắc chắn lấy được bi đỏ.\\
\textbf{Đáp án của câu hỏi là B}
}
\end{ex}

\begin{ex} %[2X13N66]
Trong một hộp chỉ có các viên bi màu đỏ. Nhắm mắt lấy ra 1 viên bi, sự kiện "Lấy được viên bi màu xanh" là
\choice
{\True Không thể}
{Chắc chắn}
{Có thể}
{Có thể hoặc không}
\loigiai{
Không có bi xanh nên không thể lấy được bi xanh.\\
\textbf{Đáp án của câu hỏi là A}
}
\end{ex}

\begin{ex} %[2X13H66]
Trong một hộp có cả bi đỏ và bi xanh. Nhắm mắt lấy ra 1 viên bi, sự kiện "Lấy được viên bi màu đỏ" là
\choice
{Chắc chắn}
{Không thể}
{Không bao giờ}
{\True Có thể}
\loigiai{
Có cả bi đỏ và bi xanh nên có thể lấy được bi đỏ.\\
\textbf{Đáp án của câu hỏi là D}
}
\end{ex}

\begin{ex} %[2X13H66]
Gieo một con xúc xắc 6 mặt (có các số chấm từ 1 đến 6). Sự kiện "Xuất hiện mặt có 7 chấm" là
\choice
{Chắc chắn}
{\True Không thể}
{Có thể}
{Có khả năng cao}
\loigiai{
Xúc xắc chỉ có từ 1 đến 6 chấm nên không thể có 7 chấm.\\
\textbf{Đáp án của câu hỏi là B}
}
\end{ex}

\begin{ex} %[2X13H66]
Gieo một con xúc xắc 6 mặt. Sự kiện "Xuất hiện mặt có số chấm bé hơn 7" là
\choice
{\True Chắc chắn}
{Không thể}
{Có thể}
{Chưa biết}
\loigiai{
Mọi mặt xúc xắc đều có số chấm từ 1 đến 6 (bé hơn 7) nên là chắc chắn.\\
\textbf{Đáp án của câu hỏi là A}
}
\end{ex}

\begin{ex} %[2X13H66]
Gieo một con xúc xắc 6 mặt. Sự kiện "Xuất hiện mặt có 3 chấm" là
\choice
{Chắc chắn}
{Không thể}
{\True Có thể}
{Hoàn toàn không}
\loigiai{
Có thể xuất hiện mặt 3 chấm.\\
\textbf{Đáp án của câu hỏi là C}
}
\end{ex}

\begin{ex} %[2X13V66]
Một túi đựng 3 quả bóng xanh và 2 quả bóng đỏ. Nam nhắm mắt lấy ra cùng lúc 3 quả bóng. Khẳng định nào sau đây là ĐÚNG?
\choice
{\True Có thể lấy được 3 quả bóng xanh}
{Chắc chắn lấy được 3 quả bóng đỏ}
{Không thể lấy được quả bóng đỏ nào}
{Chắc chắn lấy được 3 quả bóng xanh}
\loigiai{
Vì có 3 bóng xanh nên có thể lấy được 3 bóng xanh.\\
\textbf{Đáp án của câu hỏi là A}
}
\end{ex}

\begin{ex} %[2X13V66]
Một túi đựng 5 cái kẹo dâu. An lấy ngẫu nhiên 2 cái kẹo. Sự kiện "An lấy được 2 cái kẹo dâu" là
\choice
{Có thể}
{Không thể}
{\True Chắc chắn}
{Chắc chắn không xảy ra}
\loigiai{
Tất cả kẹo đều vị dâu nên chắc chắn lấy được 2 kẹo dâu.\\
\textbf{Đáp án của câu hỏi là C}
}
\end{ex}

\begin{ex} %[2X13V66]
Hôm nay là thứ Năm. Ngày mai là thứ Sáu là sự kiện
\choice
{Có thể}
{Không thể}
{Tùy từng tuần}
{\True Chắc chắn}
\loigiai{
Thứ Năm luôn đến thứ Sáu nên là chắc chắn.\\
\textbf{Đáp án của câu hỏi là D}
}
\end{ex}

\begin{ex} %[2X13V66]
Ngày mai trời sẽ có mưa là sự kiện
\choice
{Chắc chắn}
{Không thể}
{\True Có thể}
{Chắc chắn không}
\loigiai{
Thời tiết ngày mai là sự kiện có thể xảy ra.\\
\textbf{Đáp án của câu hỏi là C}
}
\end{ex}

\begin{ex} %[2X13V66]
Trong một hộp có 4 chiếc bút chì đen. Hà lấy ra 1 chiếc bút. Khẳng định nào sau đây là KHÔNG ĐÚNG?
\choice
{Chắc chắn lấy được bút chì đen}
{\True Có thể lấy được bút chì màu đỏ}
{Không thể lấy được bút chì màu đỏ}
{Chắc chắn chiếc bút lấy ra là màu đen}
\loigiai{
Không có bút đỏ nên không thể nói "có thể lấy được bút đỏ".\\
\textbf{Đáp án của câu hỏi là B}
}
\end{ex}

\begin{ex} %[2X13C66]
Trong hộp có 3 viên bi đỏ và 3 viên bi xanh. Lấy ra ngẫu nhiên 4 viên bi. Khẳng định nào sau đây là CHẮC CHẮN?
\choice
{Chắc chắn lấy được 4 viên bi đỏ}
{Chắc chắn lấy được 4 viên bi xanh}
{Chắc chắn lấy được 3 viên bi vàng}
{\True Chắc chắn trong 4 viên lấy ra có ít nhất 1 viên bi đỏ và 1 viên bi xanh}
\loigiai{
Vì chỉ có 3 bi mỗi màu, lấy 4 viên thì chắc chắn phải có cả 2 màu đỏ và xanh.\\
\textbf{Đáp án của câu hỏi là D}
}
\end{ex}

\textbf{Phần 2. Trắc nghiệm đúng sai}
\setcounterpageref{ex}{0}

\begin{ex} %[2X13H66]
Trong hộp có 5 viên bi vàng và 3 viên bi xanh. Xét tính đúng sai của các sự kiện khi lấy ra 1 viên bi:
\choiceTF
{\True Có thể lấy được viên bi màu vàng.}
{\True Có thể lấy được viên bi màu xanh.}
{Chắc chắn lấy được viên bi màu vàng.}
{\True Không thể lấy được viên bi màu đỏ.}
\loigiai{
\begin{itemize}
    \item Có thể lấy bi vàng: Đúng.
    \item Có thể lấy bi xanh: Đúng.
    \item Không chắc chắn vì có thể ra bi xanh. Ý 3 Sai, Ý 4 Đúng.
\end{itemize}
\begin{itemize}
    \item a) \textbf{Đúng}.
    \item b) \textbf{Đúng}.
    \item c) \textbf{Sai}.
    \item d) \textbf{Đúng}.
\end{itemize}
}
\end{ex}

\begin{ex} %[2X13V66]
Quay một vòng quay có 2 màu: Đỏ và Xanh. Xét tính đúng sai:
\choiceTF
{\True Kim quay có thể dừng ở phần màu đỏ.}
{\True Kim quay có thể dừng ở phần màu xanh.}
{Kim quay chắc chắn dừng ở phần màu đỏ.}
{\True Kim quay không thể dừng ở phần màu vàng (vì không có màu vàng).}
\loigiai{
\begin{itemize}
    \item Có thể dừng ở đỏ: Đúng.
    \item Có thể dừng ở xanh: Đúng.
    \item Không chắc chắn đỏ (vì có thể vào xanh). Ý 3 Sai, Ý 4 Đúng.
\end{itemize}
\begin{itemize}
    \item a) \textbf{Đúng}.
    \item b) \textbf{Đúng}.
    \item c) \textbf{Sai}.
    \item d) \textbf{Đúng}.
\end{itemize}
}
\end{ex}

\textbf{Phần 3. Trả lời ngắn}
\setcounterpageref{ex}{0}

\begin{ex} %[2X13H66]
Trong túi có 5 quả bóng đỏ. Lấy ra 1 quả bóng, sự kiện "Lấy được bóng đỏ" là gì? (Viết: chắc chắn, có thể, hoặc không thể)
\shortans{chắc chắn}
\loigiai{
\textbf{Đáp số:} chắc chắn\\
Chắc chắn.
}
\end{ex}

\begin{ex} %[2X13H66]
Trong túi có 5 quả bóng đỏ. Lấy ra 1 quả bóng, sự kiện "Lấy được bóng vàng" là gì? (Viết: chắc chắn, có thể, hoặc không thể)
\shortans{không thể}
\loigiai{
\textbf{Đáp số:} không thể\\
Không thể.
}
\end{ex}

\begin{ex} %[2X13V66]
Trong túi có 3 bóng đỏ và 3 bóng xanh. Lấy ra 1 quả bóng, sự kiện "Lấy được bóng đỏ" là gì? (Viết: chắc chắn, có thể, hoặc không thể)
\shortans{có thể}
\loigiai{
\textbf{Đáp số:} có thể\\
Có thể.
}
\end{ex}

\begin{ex} %[2X13V66]
Một con cá biết bay trên trời là sự kiện gì? (Viết: chắc chắn, có thể, hoặc không thể)
\shortans{không thể}
\loigiai{
\textbf{Đáp số:} không thể\\
Không thể.
}
\end{ex}

\begin{ex} %[2X13V66]
Mặt trời mọc ở hướng Đông là sự kiện gì? (Viết: chắc chắn, có thể, hoặc không thể)
\shortans{chắc chắn}
\loigiai{
\textbf{Đáp số:} chắc chắn\\
Chắc chắn.
}
\end{ex}

\begin{ex} %[2X13C66]
Tung một đồng xu có 2 mặt Sấp và Ngửa. Sự kiện "Đồng xu rơi xuống mặt Sấp" là gì? (Viết: chắc chắn, có thể, hoặc không thể)
\shortans{có thể}
\loigiai{
\textbf{Đáp số:} có thể\\
Có thể.
}
\end{ex}
